\chapter{Technical Implementation Details}
\label{appendix:technical}

\section{Technology Stack}

The Pulse prototype is implemented using modern web technologies optimized for serverless deployment and rapid iteration:

\subsection{Backend Framework}

\begin{itemize}
    \item \textbf{Next.js 15}: React-based framework with App Router for server-side rendering and API routes
    \item \textbf{Node.js Runtime}: JavaScript/TypeScript execution environment
    \item \textbf{Vercel}: Serverless deployment platform with edge functions
\end{itemize}

\subsection{Database and Storage}

\begin{itemize}
    \item \textbf{PostgreSQL}: Primary relational database for structured data
    \begin{itemize}
        \item Notes and folders
        \item User accounts and authentication sessions
        \item Episodic memory logs
        \item Task queue (todos)
        \item Agentic link configurations
    \end{itemize}

    \item \textbf{Vector Database}: Specialized storage for embedding-based retrieval
    \begin{itemize}
        \item Implementation: pgvector extension or dedicated vector DB (Pinecone/Weaviate)
        \item Stores: Conversation embeddings with relationship metadata
        \item Query pattern: Hybrid search (vector similarity + metadata filtering)
    \end{itemize}

    \item \textbf{Redis}: In-memory cache for session state and idempotency tokens
\end{itemize}

\subsection{AI and Language Models}

\begin{itemize}
    \item \textbf{Vercel AI SDK}: Streaming LLM responses with tool use orchestration
    \item \textbf{Model Providers}: Claude (Anthropic), GPT-4 (OpenAI), Gemini (Google)
    \item \textbf{Embedding Models}: text-embedding-3-large (OpenAI) for semantic search
\end{itemize}

\subsection{External Integrations}

\begin{itemize}
    \item \textbf{Model Context Protocol (MCP)}: Standardized interface for external tools
    \item \textbf{NextAuth.js}: OAuth 2.0 authentication for third-party services
    \item \textbf{Integrated Services}:
    \begin{itemize}
        \item Google Calendar (calendar access)
        \item Gmail (email retrieval)
        \item Microsoft Graph API (calendar/email for Microsoft users)
        \item Notion (knowledge base integration)
    \end{itemize}
\end{itemize}

\section{API Architecture}

\subsection{Primary Agent Route}

\textbf{Endpoint}: \texttt{POST /api/chat/agent-v04}

\textbf{Request Body}:
\begin{verbatim}
{
  "message": "What meetings do I have this week?",
  "conversationId": "conv_abc123",
  "safe_mode": false
}
\end{verbatim}

\textbf{Response}: Server-Sent Events (SSE) stream
\begin{verbatim}
data: {"type": "text", "content": "Let me check your calendar..."}
data: {"type": "tool_call", "name": "check_calendar", "args": {...}}
data: {"type": "tool_result", "result": [...]}
data: {"type": "text", "content": "You have 3 meetings this week..."}
data: {"type": "done"}
\end{verbatim}

\subsection{Public API Route}

\textbf{Endpoint}: \texttt{POST /api/v1/chat}

\textbf{Authentication}: API key in \texttt{Authorization} header

\textbf{Purpose}: Programmatic access for external integrations (mobile apps, CLI tools)

\subsection{Guest Interaction Route}

\textbf{Endpoint}: \texttt{POST /api/shared/[token]/chat}

\textbf{Flow}:
\begin{enumerate}
    \item Extract \texttt{token} from URL
    \item Query \texttt{agenticLinks} table for Context Cell configuration
    \item Validate token is active and not expired
    \item Mount Context Cell (filter tools, retrieve allowed note IDs)
    \item Load Relationship Shard for this guest
    \item Execute agent with restricted context
    \item Write response to guest-specific memory shard
\end{enumerate}

\section{Tool System Architecture}

\subsection{Tool Definition Schema}

Tools are defined using the OpenAI function calling schema:

\begin{verbatim}
{
  "name": "read_note",
  "description": "Retrieves the content of a specific note by ID",
  "parameters": {
    "type": "object",
    "properties": {
      "id": {
        "type": "number",
        "description": "The unique identifier of the note"
      }
    },
    "required": ["id"]
  }
}
\end{verbatim}

\subsection{Tool Categories}

\textbf{Internal Tools} (implemented directly in Pulse):
\begin{itemize}
    \item \texttt{read\_note}: Retrieve note content
    \item \texttt{search\_notes}: Keyword search across notes
    \item \texttt{create\_todo}: Add task to queue
    \item \texttt{check\_calendar}: Query calendar events
    \item \texttt{search\_email}: Retrieve email threads
\end{itemize}

\textbf{External Tools} (via MCP):
\begin{itemize}
    \item \texttt{google\_calendar\_*}: Calendar operations
    \item \texttt{gmail\_*}: Email operations
    \item \texttt{notion\_*}: Knowledge base queries
    \item \texttt{web\_search}: Internet search
\end{itemize}

\subsection{PDP Validation Logic}

\begin{verbatim}
function validateToolCall(toolCall, contextCell) {
  // Check if tool is in allowed list
  if (!contextCell.allowedTools.includes(toolCall.name)) {
    return { allowed: false, reason: "Tool not in allowedTools" };
  }

  // Validate arguments based on tool type
  if (toolCall.name === "read_note") {
    const noteId = toolCall.args.id;
    if (!contextCell.allowedNoteIds.includes(noteId)) {
      return { allowed: false, reason: `Note ${noteId} not authorized` };
    }
  }

  if (toolCall.name === "check_calendar") {
    if (contextCell.calendarAccess === "none") {
      return { allowed: false, reason: "No calendar access granted" };
    }
  }

  return { allowed: true };
}
\end{verbatim}

\section{Memory Implementation}

\subsection{Database Schema}

\textbf{Episodic Memory Table}:
\begin{verbatim}
CREATE TABLE episodic_memory (
  id SERIAL PRIMARY KEY,
  user_id INTEGER NOT NULL,
  relationship_id VARCHAR(255),  -- NULL for self-state
  conversation_id VARCHAR(255),
  role VARCHAR(50),  -- 'user' or 'assistant'
  content TEXT,
  embedding VECTOR(1536),  -- pgvector type
  timestamp TIMESTAMP DEFAULT NOW(),
  metadata JSONB
);

CREATE INDEX idx_relationship ON episodic_memory (relationship_id);
CREATE INDEX idx_embedding ON episodic_memory
  USING ivfflat (embedding vector_cosine_ops);
\end{verbatim}

\subsection{Retrieval Query}

\begin{verbatim}
SELECT content, timestamp
FROM episodic_memory
WHERE user_id = $1
  AND relationship_id = $2  -- Critical filter
ORDER BY embedding <=> $3  -- Vector similarity
LIMIT 10;
\end{verbatim}

\textbf{Key Properties}:
\begin{itemize}
    \item The \texttt{WHERE relationship\_id = \$2} clause enforces shard isolation
    \item Even if the LLM hallucinates a different relationship ID, the query is constructed server-side
    \item Zero-result guarantee if relationship ID is missing
\end{itemize}

\section{Heartbeat Engine}

\subsection{CRON Configuration}

\textbf{Trigger}: Every 15 minutes (configurable)

\textbf{Endpoint}: \texttt{POST /api/cron/heartbeat}

\textbf{Authentication}: CRON secret header

\subsection{Execution Flow}

\begin{verbatim}
async function heartbeat() {
  // 1. Fetch active users (opt-in required)
  const users = await getHeartbeatEnabledUsers();

  for (const user of users) {
    // 2. Query state deltas since last heartbeat
    const deltas = await fetchStateDeltas(user.id, lastHeartbeat);

    // 3. Invoke headless LLM to analyze deltas
    const analysis = await analyzeDelta(deltas, user.context);

    // 4. Generate actionable intents
    if (analysis.actionable) {
      await createTodo({
        user_id: user.id,
        title: analysis.intent,
        priority: analysis.urgency,
        source: 'heartbeat'
      });
    }

    // 5. Proactive notifications
    if (analysis.requiresNotification) {
      await notifyExternalParty(analysis.relationship, analysis.message);
    }
  }
}
\end{verbatim}

\section{Deployment and Observability}

\subsection{Infrastructure}

\begin{itemize}
    \item \textbf{Hosting}: Vercel (serverless functions with edge caching)
    \item \textbf{Database}: Supabase (managed PostgreSQL with pgvector)
    \item \textbf{Monitoring}: Vercel Analytics + custom logging
    \item \textbf{Error Tracking}: Sentry
\end{itemize}

\subsection{Logging Strategy}

\textbf{Log Levels}:
\begin{itemize}
    \item \texttt{INFO}: Normal operations (tool calls, memory retrievals)
    \item \texttt{WARN}: Escalations, denied tool calls
    \item \texttt{ERROR}: Exceptions, database failures
\end{itemize}

\textbf{Trace IDs}: Each request generates a unique \texttt{requestId} propagated through all function calls for end-to-end tracing.

\textbf{Audit Logs}: Security-critical events (Context Cell violations, escalations) written to dedicated audit table with immutable timestamps.

\section{Performance Optimizations}

\subsection{Caching Strategies}

\begin{itemize}
    \item \textbf{MCP Tool Discovery}: Cached per-user for 5 minutes (reduces server roundtrips)
    \item \textbf{Note Content}: Redis cache with 60-second TTL (reduces database load)
    \item \textbf{Embeddings}: Pre-computed and stored (avoids re-embedding on every query)
\end{itemize}

\subsection{Parallel Execution}

\begin{itemize}
    \item Tool calls with no dependencies executed in parallel
    \item Memory retrieval (L2 vector search) and tool mounting (Context Cell construction) run concurrently
    \item Heartbeat engine processes users in batches of 10 (rate-limited to avoid LLM quota exhaustion)
\end{itemize}

\section{Security Hardening}

\subsection{Authentication Layers}

\begin{enumerate}
    \item \textbf{User Authentication}: NextAuth session cookies
    \item \textbf{API Key Authentication}: HMAC-SHA256 signed tokens for programmatic access
    \item \textbf{Guest Token Authentication}: Time-limited, single-use tokens for external access
    \item \textbf{CRON Secret}: Environment variable verified for background job endpoints
\end{enumerate}

\subsection{Input Validation}

\begin{itemize}
    \item All user inputs sanitized against SQL injection
    \item Tool arguments validated against JSON schema before execution
    \item Note IDs validated as integers to prevent path traversal
    \item Relationship IDs validated against allowlist to prevent cross-user access
\end{itemize}

\subsection{Rate Limiting}

\begin{itemize}
    \item \textbf{Authenticated Users}: 100 requests/minute
    \item \textbf{Guest Links}: 20 requests/minute per token
    \item \textbf{Heartbeat}: 1 execution per user per 15 minutes
\end{itemize}
